\section{Namensdienst}
\label{sec:namensdienst-doku}

% ──────────────────────────────────────────────────────────────────────────────
\subsection{Konzept und Aufgabe}

Der Namensdienst (engl.\ \emph{Naming Service} oder \emph{Service Discovery})
ist eine eigenständige Systemkomponente in TaskGrid+{}.
Seine Kernaufgabe besteht darin, alle aktiven Worker zu verwalten und dem
Dispatcher auf Anfrage mitzuteilen, welche Worker für einen bestimmten
Task-Typ momentan verfügbar sind.

\noindent Ohne den Namensdienst müsste der Dispatcher Worker-Adressen statisch
kennen -- das würde dynamisches Hinzufügen, Entfernen oder Skalieren von
Workern zur Laufzeit ausschließen.
Durch die Entkopplung gilt stattdessen:

\begin{itemize}
  \item \textbf{Worker melden sich selbst an} (Push-Modell):
        Jeder Worker registriert sich beim Start und bestätigt seine
        Lebendigkeit durch regelmäßige Heartbeats.
  \item \textbf{Dispatcher fragt bei Bedarf ab} (Pull-Modell):
        Bevor ein Task dispatcht wird, fragt der Dispatcher den Namensdienst
        per \texttt{LookupWorker} nach passenden Workern.
  \item \textbf{Kein statisches Wissen im Dispatcher}:
        Worker-Adressen tauchen in keiner Konfigurationsdatei des Dispatchers
        auf; sie existieren ausschließlich in der In-Memory-Registry des
        Namensdiensts.
\end{itemize}

\noindent Die Abgrenzung gegenüber Docker-DNS ist in Abschnitt~\ref{sec:nd-docker-dns}
beschrieben.

% ──────────────────────────────────────────────────────────────────────────────
\subsection{Datenmodell -- Wie werden Worker gespeichert?}
\label{sec:nd-datenmodell}

Der Namensdienst hält alle Registrierungen \textbf{ausschließlich im
Arbeitsspeicher} (Python-Liste).
Eine Persistenzschicht (Datenbank, Datei, Volume) ist bewusst nicht
vorgesehen (§4.1 der Aufgabenstellung):
Stale-Einträge werden durch den Heartbeat-Timeout automatisch entfernt,
und neue Worker registrieren sich beim eigenen Container-Start selbst neu.

\noindent Jeder registrierte Worker wird durch eine \texttt{Worker}-Instanz
(in \texttt{src/namensdienst/worker.py}) repräsentiert:

\begin{table}[H]
\centering
\begin{tabular}{m{3.5cm} m{2.4cm} m{7.5cm}}
\toprule
\textbf{Attribut} & \textbf{Typ} & \textbf{Bedeutung} \\
\midrule
\texttt{worker\_id}    & \texttt{str}   & Eindeutige ID des Workers
                                          (Format: \texttt{worker-<typ>-<hostname>}) \\
\texttt{type}          & \texttt{str}   & Primärer Task-Typ dieses Eintrags
                                          (pro Typ ein eigener Eintrag) \\
\texttt{task\_types}   & \texttt{list}  & Alle Task-Typen, die der Worker unterstützt \\
\texttt{address}       & \texttt{str}   & Docker-Compose-Service-Name
                                          (z.\,B. \texttt{worker-sum}) \\
\texttt{port}          & \texttt{int}   & gRPC-Port des Workers im Container \\
\texttt{status}        & \texttt{str}   & Zustand: \texttt{ACTIVE}, \texttt{UNHEALTHY}
                                          oder \texttt{OFFLINE} \\
\texttt{lastHeartbeat} & \texttt{float} & Unix-Timestamp des letzten Heartbeats \\
\texttt{currentLoad}   & \texttt{int}   & Aktuelle Anzahl gleichzeitig bearbeiteter Tasks \\
\texttt{id}            & \texttt{int}   & Interner Auto-Inkrement-Zähler des Namensdiensts \\
\bottomrule
\end{tabular}
\caption{Felder der \texttt{Worker}-Klasse im Namensdienst}
\end{table}

\noindent Unterstützt ein Worker mehrere Task-Typen (z.\,B. \texttt{sum} und
\texttt{reverse}), legt der Namensdienst \textbf{pro Typ einen eigenen Eintrag}
an -- alle mit derselben \texttt{worker\_id}.
Dadurch ist \texttt{LookupWorker} eine einfache Filteroperation
auf \texttt{worker.type == task\_type}.

Der Zähler \texttt{idcount} in der \texttt{Namensdienst}-Klasse
stellt sicher, dass jeder interne Eintrag eine monoton steigende
numerische ID erhält, auch wenn Workers sich mehrfach registrieren.

% ──────────────────────────────────────────────────────────────────────────────
\subsection{Registrierungsablauf}
\label{sec:nd-registrierung-ablauf}

Ein Worker sendet beim Start einen \texttt{RegisterWorker}-gRPC-Aufruf
an den Namensdienst.
Auf Server-Seite läuft dabei folgendes ab:

\begin{enumerate}
  \item \textbf{Validierung}: Ist \texttt{worker\_id} leer, wird der Aufruf
        mit \texttt{INVALID\_ARGUMENT} abgelehnt.
  \item \textbf{Idempotenz}: Alle bestehenden Einträge mit gleicher
        \texttt{worker\_id}, \texttt{address} und \texttt{port} werden
        vor dem Einfügen entfernt.
        Damit ist eine erneute Registrierung (z.\,B. nach Neustart des Workers)
        sicher möglich, ohne Duplikate zu erzeugen.
  \item \textbf{Einträge anlegen}: Für jeden unterstützten Task-Typ wird
        ein separater \texttt{Worker}-Eintrag in \texttt{self.workers}
        eingefügt.
        \texttt{lastHeartbeat} wird auf die aktuelle Zeit gesetzt,
        \texttt{status} auf den übergebenen Wert (typisch \texttt{ACTIVE}).
  \item \textbf{Antwort}: Der Namensdienst antwortet mit
        \texttt{Ack(success=True)}.
\end{enumerate}

\begin{lstlisting}[language=Python,caption={Registrierungslogik -- Namensdienst-seitig (nameservice.py)}]
def RegisterWorker(self, request, context):
    payload = request.payload
    # Schritt 1: Validierung
    if not payload.worker_id:
        context.set_code(grpc.StatusCode.INVALID_ARGUMENT)
        return taskgrid_pb2.Ack(...)
    # Schritt 2: Duplikate entfernen (Idempotenz)
    self.workers = [
        w for w in self.workers
        if not (w.worker_id == payload.worker_id
                and w.address == payload.address
                and w.port    == payload.port)
    ]
    # Schritt 3: Einen Eintrag pro Task-Typ anlegen
    for task_type in task_types:
        worker = Worker(task_type, payload.address, payload.port,
                        self.idcount, self,
                        worker_id=payload.worker_id,
                        task_types=task_types)
        worker.lastHeartbeat = time.time()
        worker.status = payload.status or "ACTIVE"
        self.workers.append(worker)
        self.idcount += 1
\end{lstlisting}

% ──────────────────────────────────────────────────────────────────────────────
\subsection{Lookup-Ablauf}
\label{sec:nd-lookup-ablauf}

Der Dispatcher ruft \texttt{LookupWorker} auf, bevor er einen Task
an einen Worker sendet.
Der Aufruf enthält im Payload lediglich den gesuchten \texttt{task\_type}.

\noindent Der Namensdienst filtert intern alle \texttt{Worker}-Einträge:

\begin{lstlisting}[language=Python,caption={LookupWorker-Implementierung (nameservice.py)}]
def LookupWorker(self, request, context):
    task_type = request.payload.task_type
    result = [
        w for w in self.workers
        if w.type == task_type
        and w.status not in ("UNHEALTHY", "OFFLINE")
    ]
    return taskgrid_pb2.LookupResponse(
        payload=taskgrid_pb2.LookupResponse.Payload(
            found=bool(result),
            workers=[
                taskgrid_pb2.LookupResponse.WorkerPayload(
                    worker_id=w.worker_id,
                    task_types=w.task_types,
                    address=w.address,
                    port=int(w.port),
                    status=w.status,
                    current_load=w.currentLoad,
                )
                for w in result
            ],
            message="Workers found" if result else "No workers found",
        )
    )
\end{lstlisting}

\noindent Worker mit Status \texttt{UNHEALTHY} oder \texttt{OFFLINE} werden
\textbf{nicht} zurückgegeben.
Gibt es keine passenden Worker, enthält die Antwort
\texttt{found=False} und eine leere Worker-Liste.
Der Dispatcher behandelt eine leere Liste als
"`kein Worker verfügbar"' und legt den Task nach
\texttt{DISPATCH\_NO\_WORKER\_RETRY\_SECONDS} (Standard: 2\,s) erneut
in die Queue.

% ──────────────────────────────────────────────────────────────────────────────
\subsection{Abgrenzung zu Docker-DNS}
\label{sec:nd-docker-dns}

Docker Compose bietet mit seinem eingebauten DNS-Dienst eine Möglichkeit,
Container-Namen in IP-Adressen aufzulösen.
Der eigene Namensdienst in TaskGrid+ erfüllt jedoch eine andere,
erweiterte Funktion:

\begin{table}[H]
\centering
\begin{tabular}{m{4.5cm} m{4.5cm} m{4.5cm}}
\toprule
\textbf{Eigenschaft} & \textbf{Docker-DNS} & \textbf{TaskGrid+-Namensdienst} \\
\midrule
Auflösung           & Hostname → IP       & Worker-Typ → Adresse(n), Port, Status \\
Gesundheitszustand  & Nicht bekannt       & Heartbeat-Status (\texttt{ACTIVE}/\texttt{UNHEALTHY}) \\
Lastinformation     & Nicht verfügbar     & \texttt{current\_load} pro Worker \\
Mehrere Instanzen   & Round-Robin-DNS     & Explizite Liste aller Instanzen \\
Registrierung       & Automatisch         & Explizit durch Worker selbst \\
Fehler bei Ausfall  & DNS-Fehler          & Leere Liste, kein Absturz \\
\bottomrule
\end{tabular}
\caption{Vergleich: Docker-DNS vs. TaskGrid+-Namensdienst}
\end{table}

\noindent Docker-DNS kann zwar mehrere IP-Adressen für einen Service-Namen
zurückgeben (Round-Robin), kennt aber weder den Gesundheitszustand
einzelner Container noch deren aktuellen Auslastungsgrad.
Der eigene Namensdienst ermöglicht es, ausgefallene oder überlastete
Worker aus der Auswahl auszuschließen, bevor der Dispatcher überhaupt
versucht, einen Task zu senden.

% ──────────────────────────────────────────────────────────────────────────────
\subsection{Worker-Registrierung und Heartbeats}
\label{sec:nd-heartbeat}

\subsubsection*{Registrierungsablauf beim Container-Start}

Beim Start eines Worker-Containers durchläuft \texttt{serve()} in
\texttt{src/worker/worker.py} folgende Schritte:

\begin{enumerate}
  \item Signal-Handler für \texttt{SIGTERM} und \texttt{SIGINT} registrieren
        (Graceful Shutdown).
  \item \texttt{register\_worker()} aufrufen -- bis zu 5 Versuche mit
        exponentiellem Backoff (1\,s, 2\,s, 4\,s, 8\,s, 16\,s).
  \item Heartbeat-Thread starten (\texttt{send\_heartbeat\_loop},
        Daemon-Thread, läuft parallel zur Task-Verarbeitung).
  \item gRPC-Server starten und auf eingehende \texttt{ExecuteTask}-Aufrufe
        warten.
\end{enumerate}

\noindent Die Identität des Workers wird durch drei Umgebungsvariablen
festgelegt, die im \texttt{entrypoint.sh} gesetzt werden:

\begin{lstlisting}[language=bash,caption={Worker-Identität (entrypoint.sh)}]
# WORKER_ID: eindeutig pro Container-Instanz
if [ -z "${WORKER_ID}" ]; then
    export WORKER_ID="worker-${TASK_TYPES}-$(hostname)"
fi
# NAMING_SERVICE_ADDRESS: NAMENSDIENST_ADDRESS hat Vorrang
export NAMING_SERVICE_ADDRESS=\
    "${NAMENSDIENST_ADDRESS:-${NAMING_SERVICE_ADDRESS:-namensdienst:50052}}"
\end{lstlisting}

\noindent \texttt{WORKER\_HOST} muss dem Docker-Compose-Service-Namen entsprechen
(z.\,B. \texttt{worker-sum}), damit Docker-DNS den Namen zur Container-IP
auflösen kann, wenn der Dispatcher den Worker über die vom Namensdienst
zurückgegebene Adresse kontaktiert.

\subsubsection*{Heartbeat-Intervall und konfigurierbare Parameter}

Der Heartbeat-Loop sendet in fester Periode einen
\texttt{SendHeartbeat}-gRPC-Aufruf an den Namensdienst.
Dabei werden \texttt{worker\_id} und \texttt{current\_load}
übertragen, sodass der Namensdienst die aktuelle Auslastung stets kennt.

\begin{lstlisting}[language=Python,caption={Heartbeat-Loop (worker.py)}]
def send_heartbeat_loop():
    while True:
        stub.SendHeartbeat(
            taskgrid_pb2.HeartbeatRequest(
                message_type="HEARTBEAT",
                sender=WORKER_ID,
                timestamp=int(time.time()),
                request_id=str(uuid.uuid4()),
                payload=taskgrid_pb2.HeartbeatRequest.Payload(
                    worker_id=WORKER_ID,
                    current_load=get_current_load()
                )
            )
        )
        time.sleep(HEARTBEAT_INTERVAL)
\end{lstlisting}

\noindent Alle zeitrelevanten Parameter sind über Umgebungsvariablen konfigurierbar:

\begin{table}[H]
\centering
\begin{tabular}{m{6.5cm} m{2cm} m{5cm}}
\toprule
\textbf{Umgebungsvariable} & \textbf{Default} & \textbf{Wirkung} \\
\midrule
\texttt{HEARTBEAT\_INTERVAL\_SECONDS}
  & 5\,s
  & Sendeintervall des Workers \\
\texttt{NAMESERVICE\_UNHEALTHY\_SECS}
  & 10\,s
  & Namensdienst: Schwelle für \texttt{UNHEALTHY} \\
\texttt{NAMESERVICE\_OFFLINE\_MULT}
  & 2
  & Multiplikator für \texttt{OFFLINE}-Schwelle
    ($= \texttt{UNHEALTHY} \times \texttt{MULT}$) \\
\bottomrule
\end{tabular}
\caption{Konfigurierbare Heartbeat-Parameter}
\end{table}

\noindent Mit den Standardwerten ergibt sich:
UNHEALTHY nach 10\,s ohne Heartbeat,
OFFLINE nach 20\,s ohne Heartbeat
(entspricht \texttt{NAMESERVICE\_UNHEALTHY\_SECS} $\times$
\texttt{NAMESERVICE\_OFFLINE\_MULT}).

\subsubsection*{Inaktivitätserkennung -- wann wird ein Worker UNHEALTHY bzw.\ OFFLINE?}

Im Namensdienst läuft ein Hintergrund-Thread (\texttt{startLoop}), der
jede Sekunde alle registrierten Worker überprüft.
Die Entscheidung basiert auf der vergangenen Zeit seit dem letzten
Heartbeat:

\begin{lstlisting}[language=Python,caption={Inaktivitätserkennung (nameservice.py)}]
def loop():
    while self.running:
        now = time.time()
        for worker in self.workers[:]:          # Kopie: sicheres Entfernen
            elapsed = now - worker.lastHeartbeat
            if elapsed >= x:                    # x = NAMESERVICE_UNHEALTHY_SECS
                worker.status = "UNHEALTHY"
            if elapsed >= x * y:                # y = NAMESERVICE_OFFLINE_MULT
                worker.status = "OFFLINE"
                self.workers.remove(worker)
        time.sleep(1)
\end{lstlisting}

\noindent Die Zustandsübergänge lauten damit:

\begin{enumerate}
  \item \textbf{ACTIVE} $\rightarrow$ \textbf{UNHEALTHY}:
        Der letzte Heartbeat liegt länger als
        \texttt{NAMESERVICE\_UNHEALTHY\_SECS} (Standard: 10\,s) zurück.
        Der Worker wird weiterhin in \texttt{self.workers} geführt,
        aber von \texttt{LookupWorker} \textbf{nicht mehr zurückgegeben}.
  \item \textbf{UNHEALTHY} $\rightarrow$ \textbf{OFFLINE}:
        Der letzte Heartbeat liegt länger als
        $\texttt{UNHEALTHY\_SECS} \times \texttt{OFFLINE\_MULT}$
        (Standard: 20\,s) zurück.
        Der Worker wird \textbf{aus der Liste entfernt} und gilt als
        endgültig nicht mehr verfügbar.
\end{enumerate}

\noindent Ein als \texttt{OFFLINE} entfernter Worker kann sich anschließend
erneut registrieren (z.\,B. nach Container-Neustart) -- der
Namensdienst behandelt jede Registrierung idempotent
(siehe Abschnitt~\ref{sec:nd-registrierung-ablauf}).

\noindent Beim \textbf{Graceful Shutdown} sendet der Worker \emph{keine}
explizite Abmeldung (\texttt{DeregisterWorker}), sondern stoppt lediglich
den Heartbeat-Loop.
Damit durchläuft der Namensdienst die Zustands\-übergänge
ACTIVE $\rightarrow$ UNHEALTHY $\rightarrow$ OFFLINE selbstständig
und testbar -- ohne dass ein manueller Eingriff nötig wäre.

% ──────────────────────────────────────────────────────────────────────────────
\subsection{Sequenzdiagramm: Worker-Start, Registrierung und Heartbeat-Loop}
\label{sec:nd-sequenz}

Das folgende Diagramm zeigt den vollständigen Ablauf vom Container-Start
eines Workers bis zum laufenden Heartbeat-Betrieb.

\begin{figure}[H]
\centering
\includegraphics[width=\textwidth]{NamensdienstSequenz.drawio.png}
\caption{Sequenzdiagramm: Worker-Start -- Register -- Heartbeat-Loop}
\end{figure}

\noindent Da das Diagramm noch nicht als Bilddatei vorliegt, wird der Ablauf
hier tabellarisch beschrieben:

\begin{table}[H]
\centering
\begin{tabular}{m{0.5cm} m{2.5cm} m{2.5cm} m{8cm}}
\toprule
\textbf{\#} & \textbf{Von} & \textbf{An} & \textbf{Nachricht / Aktion} \\
\midrule
1  & Docker & Worker       & Container startet, \texttt{entrypoint.sh} setzt Umgebungsvariablen \\
2  & Worker & Namensdienst & \texttt{RegisterWorker(worker\_id, task\_types, address, port, status=ACTIVE)} \\
3  & Namensdienst & Worker & \texttt{Ack(success=True)} \\
4  & Worker & Worker       & Heartbeat-Thread starten (\texttt{send\_heartbeat\_loop}) \\
5  & Worker & Worker       & gRPC-Server starten (\texttt{WorkerService}) \\
6  & Worker & Namensdienst & \texttt{SendHeartbeat(worker\_id, current\_load)}
                             $\quad$ (alle \texttt{HEARTBEAT\_INTERVAL\_SECONDS}) \\
7  & Namensdienst & Worker & \texttt{Ack(success=True, message="Heartbeat received")} \\
8  & Namensdienst & --     & Hintergrundloop prüft \texttt{lastHeartbeat} jede Sekunde \\
9  & -- & --               & Falls kein Heartbeat $\geq$ \texttt{UNHEALTHY\_SECS}: Status $\rightarrow$ UNHEALTHY \\
10 & -- & --               & Falls kein Heartbeat $\geq$ \texttt{UNHEALTHY\_SECS} $\times$ \texttt{MULT}: Status $\rightarrow$ OFFLINE, Eintrag entfernt \\
\bottomrule
\end{tabular}
\caption{Schrittweise Beschreibung des Sequenzdiagramms}
\end{table}